\section{Register Pointer, Index, dan Instruction Pointer}
\label{sec:register-pointer}

Kelompok register berikutnya berjumlah empat buah: \reg{SP}, \reg{BP},
\reg{SI}, dan \reg{DI}. Keempatnya berukuran 16 bit dan menyimpan
\eng{offset}, yaitu jarak suatu lokasi dari awal segmen. Karena itu,
kelompok ini disebut register \eng{pointer} dan \eng{index}.
Pembahasan berikut diringkas dari \cite{gunawan2019}, Bagian 2,
hlm.~13--18, dan \cite{sto2001}, BAB IV, hlm.~17--20.

\subsection{Stack Pointer dan Base Pointer}

\eng{Stack} adalah daerah memori yang bekerja seperti tumpukan:
data yang terakhir dimasukkan adalah data yang pertama diambil
(\eng{last in, first out}). Instruksi \asm{PUSH} menaruh data ke
\eng{stack}, sedangkan \asm{POP} mengambilnya kembali.

\begin{description}
    \item[\reg{SP} (\eng{Stack Pointer})] Menunjuk puncak
          \eng{stack} yang aktif. Pasangannya adalah \reg{SS},
          sehingga alamat puncak \eng{stack} ditulis sebagai
          \reg{SS}:\reg{SP}. Nilai \reg{SP} berubah ketika
          \asm{PUSH} dan \asm{POP} dijalankan.

    \item[\reg{BP} (\eng{Base Pointer})] Menjadi titik acuan di
          dalam \eng{stack}. Pasangannya juga \reg{SS}, sehingga
          alamatnya ditulis \reg{SS}:\reg{BP}. \reg{BP} berguna
          untuk menunjuk data pada \eng{stack} tanpa terpengaruh
          perubahan \reg{SP}.
\end{description}

Karena \reg{SP} menunjukkan puncak \eng{stack}, mengubah nilainya
secara sembarangan berarti memindahkan seluruh tumpukan. Akibatnya,
instruksi \asm{RET} dapat melompat ke alamat yang salah.

\subsection{Source Index dan Destination Index}

Dua register berikutnya dipakai untuk mengakses deretan data,
misalnya string atau \eng{array}.

\begin{description}
    \item[\reg{SI} (\eng{Source Index})] Menunjuk alamat sumber pada
          operasi string. Pasangannya adalah \reg{DS}, sehingga
          alamat sumber ditulis \reg{DS}:\reg{SI}.

    \item[\reg{DI} (\eng{Destination Index})] Menunjuk alamat tujuan
          pada operasi string. Pasangannya adalah \reg{ES}, sehingga
          alamat tujuan ditulis \reg{ES}:\reg{DI}.
\end{description}

Pada operasi string, \reg{SI} dan \reg{DI} dapat berubah secara
otomatis mengikuti arah pemrosesan. Karena itu, kedua register harus
diisi dengan alamat yang benar sebelum operasi dijalankan.

\subsection{Instruction Pointer}

Register terakhir yang perlu dikenal adalah \reg{IP}
(\eng{Instruction Pointer}). Register ini menyimpan \eng{offset}
instruksi berikutnya yang akan diambil dan dieksekusi. Pasangan
\reg{CS}:\reg{IP} menentukan alamat instruksi tersebut. Nilai
\reg{IP} bertambah sendiri setiap kali instruksi diambil.

Programmer hampir tidak pernah menulis \reg{IP} secara langsung.
Perubahan \reg{IP} terjadi melalui instruksi seperti \asm{JMP},
\asm{CALL}, \asm{RET}, dan \asm{INT}.

Gambar~\ref{fig:organisasi-register} merangkum pembagian keempat
belas register 8086 ke dalam kelompoknya.

\begin{figure}[H]
    \centering
    \begin{tikzpicture}[
        reg/.style={
            draw=obeBlue,
            line width=0.7pt,
            fill=obeBlueLight,
            rounded corners=2pt,
            minimum width=1.5cm,
            minimum height=0.62cm,
            align=center,
            font=\small\ttfamily,
            inner sep=3pt
        },
        regk/.style={reg, fill=obeTealLight, draw=obeTeal},
        regs/.style={reg, fill=obeAmberLight, draw=obeAmber},
        regf/.style={reg, fill=obeGrayLight, draw=obeGray},
        judul/.style={font=\small\bfseries\sffamily, text=obeBlue},
        ket/.style={font=\scriptsize\itshape, text=obeGray,
                    align=center, text width=1.45cm}
    ]
        % --- kelompok 1: register umum ---
        \node[judul] at (0,4.6) {Register umum};
        \node[reg] (ax) at (0,4.05) {AX};
        \node[reg, right=0.25cm of ax] (bx) {BX};
        \node[reg, right=0.25cm of bx] (cx) {CX};
        \node[reg, right=0.25cm of cx] (dx) {DX};
        \node[ket, below=0.10cm of ax] {akumulator};
        \node[ket, below=0.10cm of bx] {offset};
        \node[ket, below=0.10cm of cx] {pencacah};
        \node[ket, below=0.10cm of dx] {data};

        % --- kelompok 2: register segmen ---
        \node[judul] at (0,2.75) {Register segmen};
        \node[regk] (cs) at (0,2.25) {CS};
        \node[regk, right=0.25cm of cs] (ds) {DS};
        \node[regk, right=0.25cm of ds] (ss) {SS};
        \node[regk, right=0.25cm of ss] (es) {ES};
        \node[ket, below=0.10cm of cs] {kode};
        \node[ket, below=0.10cm of ds] {data};
        \node[ket, below=0.10cm of ss] {stack};
        \node[ket, below=0.10cm of es] {tambahan};

        % --- kelompok 3: pointer dan index ---
        \node[judul] at (0,0.95) {Pointer, index, dan IP};
        \node[regs] (sp) at (0,0.45) {SP};
        \node[regs, right=0.25cm of sp] (bp) {BP};
        \node[regs, right=0.25cm of bp] (si) {SI};
        \node[regs, right=0.25cm of si] (di) {DI};
        \node[regs, right=0.25cm of di] (ip) {IP};
        \node[ket, below=0.10cm of sp] {puncak stack};
        \node[ket, below=0.10cm of bp] {acuan stack};
        \node[ket, below=0.10cm of si] {sumber};
        \node[ket, below=0.10cm of di] {tujuan};
        \node[ket, below=0.10cm of ip] {instruksi};

        % --- kelompok 4: flags ---
        \node[judul] at (0,-0.75) {Register flag};
        \node[regf, minimum width=7.4cm] (fl) at (1.0,-1.35)
            {FLAGS (16 bit: CF PF AF ZF SF TF IF DF OF)};

        % --- garis pengelompokan ---
        \begin{scope}[on background layer]
            \node[draw=obeBlue, dashed, rounded corners=4pt,
                  fit=(ax)(dx), inner sep=6pt] {};
            \node[draw=obeTeal, dashed, rounded corners=4pt,
                  fit=(cs)(es), inner sep=6pt] {};
            \node[draw=obeAmber, dashed, rounded corners=4pt,
                  fit=(sp)(ip), inner sep=6pt] {};
            \node[draw=obeGray, dashed, rounded corners=4pt,
                  fit=(fl), inner sep=6pt] {};
        \end{scope}
    \end{tikzpicture}
    \caption{Organisasi keempat belas register 8086 menurut
    kelompoknya.}
    \label{fig:organisasi-register}
\end{figure}

\subsection{Daftar Ringkas Register 8086}

Tabel~\ref{tab:daftar-register} merangkum seluruh register beserta
kelompok dan fungsi utamanya.

\begin{table}[H]
    \centering
    \caption{Daftar ringkas register 8086.}
    \label{tab:daftar-register}
    \small
    \begin{tabular}{@{}llp{2.6cm}p{6.8cm}@{}}
        \toprule
        \textbf{Register} & \textbf{Ukuran} & \textbf{Kelompok} &
        \textbf{Fungsi utama} \\
        \midrule
        \reg{AX} & 16 bit & Umum &
        Akumulator; dapat diakses sebagai \reg{AH} dan \reg{AL} \\
        \addlinespace
        \reg{BX} & 16 bit & Umum &
        Penunjuk alamat \eng{offset}; dapat diakses sebagai \reg{BH}
        dan \reg{BL} \\
        \addlinespace
        \reg{CX} & 16 bit & Umum &
        Pencacah pengulangan; dapat diakses sebagai \reg{CH} dan
        \reg{CL} \\
        \addlinespace
        \reg{DX} & 16 bit & Umum &
        Data pendukung dan sisa pembagian; dapat diakses sebagai
        \reg{DH} dan \reg{DL} \\
        \addlinespace
        \reg{CS} & 16 bit & Segmen &
        Menunjuk segmen kode; berpasangan dengan \reg{IP} \\
        \addlinespace
        \reg{DS} & 16 bit & Segmen &
        Menunjuk segmen data \\
        \addlinespace
        \reg{SS} & 16 bit & Segmen &
        Menunjuk segmen \eng{stack}; berpasangan dengan \reg{SP} dan
        \reg{BP} \\
        \addlinespace
        \reg{ES} & 16 bit & Segmen &
        Segmen tambahan; berpasangan dengan \reg{DI} pada operasi
        string \\
        \addlinespace
        \reg{SP} & 16 bit & Pointer &
        Menunjuk puncak \eng{stack} \\
        \addlinespace
        \reg{BP} & 16 bit & Pointer &
        Titik acuan data di dalam \eng{stack} \\
        \addlinespace
        \reg{SI} & 16 bit & Index &
        Indeks sumber \reg{DS}:\reg{SI} \\
        \addlinespace
        \reg{DI} & 16 bit & Index &
        Indeks tujuan \reg{ES}:\reg{DI} \\
        \addlinespace
        \reg{IP} & 16 bit & Kendali &
        \eng{Offset} instruksi berikutnya; berpasangan dengan
        \reg{CS} \\
        \addlinespace
        \reg{FLAGS} & 16 bit & Status &
        Sembilan \eng{flag} kondisi dan kendali \\
        \bottomrule
    \end{tabular}
\end{table}

\begin{infobox}
\vspace{1.1em}
Aturan praktis: register umum \reg{AX}--\reg{DX} bebas dipakai untuk
perhitungan; \reg{SI}, \reg{DI}, dan \reg{BX} dapat menjadi penunjuk
alamat selama isinya dicatat; \reg{CS}, \reg{IP}, \reg{SS}, dan
\reg{SP} hanya diubah melalui instruksi yang memang dirancang untuk
itu.
\end{infobox}
